Reputation systems are widely assumed to stabilize cooperation.
We ask whether that holds for large language model (LLM) agents in a climate-risk public-goods economy with asymmetric institutions.
Twenty-six agents interact for thirty rounds under climatic shocks, a dual-use loss-and-damage pool, democratic rule updating, and social-information channels that include peer scoring, reputation, and gossip.
Two independent replications share the full protocol and differ only by random seed.
Contribution intensity relative to wealth falls, on average, after bad-reputation and gossip events among agents in the non-enforcing institution.
Responses among enforcing-institution agents reverse sign across replications.
Average cooperation remains moderately positive, yet democratic adaptation and wealth trajectories diverge.
Social monitoring can coincide with withdrawal when agents lack peer-enforcement tools.
The finding challenges the default claim that reputation raises cooperation in LLM climate societies.
